\section{Regional clustering: methods and conclusions}
\textbf{Reproduce:} \path{Rscript scripts/table_main.R}. Helper:
\path{scripts/cluster_inference.R}. Checks:
\path{Rscript tests/check_table2_clustering.R}.
Numerical results: \path{output/tables/tableII_mod_diagnostics.csv}.
Region-level sample sizes and AR scores: \path{output/tables/tableII_mod_cluster_audit.csv}.

All original 2SLS point estimates, samples and controls are retained.
Clustering changes the inference. CR2 bias-adjusted regional variances use
coefficient-specific Satterthwaite degrees of freedom, implemented by
\texttt{clubSandwich} on equivalent \texttt{AER::ivreg} fits. Each fit is checked
against the original \texttt{felm} estimate. These are small-sample approximations,
not exact tests, and Wald inference also requires adequate identification.

\begin{center}
\begin{tabular}{lrrrr}
\toprule
 & \multicolumn{2}{c}{Whole sample} & \multicolumn{2}{c}{Excluding Amhara}\\
 & Salary & Wage & Salary & Wage\\\midrule
Estimate & .0431 & .0292 & .0423 & .0289\\
Conventional $p$ & .007 & .277 & .018 & .322\\
CR2 SE & .0183 & .0095 & .0171 & .0169\\
Satterthwaite d.f. & 1.62 & 1.95 & 1.31 & 1.44\\
CR2 $p$ & .172 & .094 & .194 & .276\\
AR CR2 $p$ & .123 & .274 & .081 & .438\\
AR sign-flip $p$ & .250 & .250 & .500 & .500\\\bottomrule
\end{tabular}
\end{center}
\textbf{Conclusion.} Salary employment no longer rejects a zero effect at 5\%.
None of the four CR2 tests rejects at 5\%; all CR2 95\% confidence intervals include
zero. Whole-sample wage is marginal at 10\% under the Wald check, while the
corresponding AR test is not. Conversely, the excluding-Amhara salary AR test is
marginal at 10\%, while its Wald test is not. Neither is stable across methods.
The evidence cannot establish a robust positive employment effect or a precise zero.

\textbf{Weak-instrument check.} The CR2 AR test applies the same adjustment to
$y$ on controls and $E_{is}$ under $H_0:\beta=0$. It avoids dividing by an estimated
first-stage coefficient, but still relies on small-sample approximations.
Clustered first-stage F statistics, d.f. and p-values are recorded in the CSV;
with 3--4 regions they should not be assessed using many-cluster weak-IV cutoffs.

\textbf{Sign-flip sensitivity.} Inspired by Wang--Zhang's wild AR score procedure,
the instrument is orthogonalized against controls separately by region to avoid
assuming the same instrument--control projection across regions. Null-restricted
scores are aggregated by region and all $2^G$ Rademacher patterns are enumerated.
This is a different set of test moments from the pooled instrument: it discards
between-region variation, while leaving the reported point estimator unchanged.
Tigray's residual instrument is zero; only three regions, or two excluding Amhara,
are informative. Tied patterns are included. Attainable minimum two-sided
p-values are .250 and .500. Additional draws would merely repeat these patterns.

The check requires independent regions, conditional instrument validity within
regions, and approximately symmetric cluster scores, supported asymptotically
by suitable within-region weak dependence. It does not accommodate arbitrary
common regional factor shocks, justify the exclusion restriction, or provide an
exact finite-sample guarantee. Policy assignments are not permuted: no known
random assignment mechanism is assumed. Webb weights or smaller administrative
clusters should not be selected simply to obtain finer or smaller p-values.

\textbf{Method sources:}
\href{https://doi.org/10.1080/07350015.2016.1247004}{Pustejovsky and Tipton (2018)},
\href{https://jepusto.github.io/clubSandwich/}{clubSandwich documentation}, and
\href{https://doi.org/10.1016/j.jeconom.2024.105727}{Wang and Zhang (2024)}.
